\chapter[12]{Weak and Strong Maximum Principles on Noncompact Manifolds}
\label{ch:noncompact-maximum-principles}

\section{Weak maximum principles for scalar heat-type equations}

\subsection{The heat equation on Euclidean space}

\begin{theorem}[WMP on $\mathbb R^n$ under exponential quadratic growth]
\label{thm:euclidean-heat-wmp-exponential-quadratic-growth}
\source{chow-et-al-ricci-flow-part-ii:page-0168}
\dcref{ch12:12.1}
\group{Weak Scalar Maximum Principles}

Let $u\in C^{2,1}(\mathbb R^n\times(0,T])\cap
C^0(\mathbb R^n\times[0,T])$ solve $u_t-\Delta u=0$.  If
\[
u(x,t)\leq A e^{B|x|^2}
\]
for finite constants $A,B$, then
\[
\sup_{\mathbb R^n\times[0,T]}u=\sup_{\mathbb R^n}u(\cdot,0).
\]
\end{theorem}

\begin{proof}
\sketch
\source{chow-et-al-ricci-flow-part-ii:page-0168,chow-et-al-ricci-flow-part-ii:page-0169}
Write $D=\sup u(\cdot,0)$, the case $D=+\infty$ being immediate.  Choose
$T'\leq T$ with $4BT'<1$ and, for $\delta,\varepsilon>0$, set
\[
v=u-\delta(T'+\varepsilon-t)^{-n/2}
 \exp\!\left(\frac{|x|^2}{4(T'+\varepsilon-t)}\right),
\qquad H=v-D-\delta t.
\]
For sufficiently small $\varepsilon$, the subtracted term grows faster than
$Ae^{B|x|^2}$, uniformly for $0\leq t\leq T'$, so $H\to-\infty$ at spatial
infinity.  If $H$ first attained the value $0$ at $(x_0,t_0)$ with $t_0>0$,
then $H_t\geq0$ and $\Delta H\leq0$ there.  Since $v_t=\Delta v$, this gives
$0\leq H_t=\Delta H-\delta<0$, a contradiction.  Thus
$u\leq D+\delta T'$ on the first interval.  Letting $\delta\downarrow0$ and
repeating the argument on consecutive intervals of length $T'$ proves the
claim.
\end{proof}

\begin{remark}[The barrier method]
\label{rem:subtract-little-of-a-lot-barrier}
\source{chow-et-al-ricci-flow-part-ii:page-0169}
\uses{thm:euclidean-heat-wmp-exponential-quadratic-growth}
\dcref{ch12:12.2}
\group{Weak Scalar Maximum Principles}


The proof subtracts a small multiple of a rapidly growing solution.  The
resulting function tends to $-\infty$ at spatial infinity and therefore has a
finite maximum point to which the second derivative test applies.
\end{remark}

\begin{theorem}[Uniqueness on $\mathbb R^n$ under exponential quadratic growth]
\label{thm:euclidean-heat-uniqueness-exponential-quadratic-growth}
\source{chow-et-al-ricci-flow-part-ii:page-0169}
\dcref{ch12:12.3}
\group{Weak Scalar Maximum Principles}

Let $u_1,u_2\in C^{2,1}(\mathbb R^n\times(0,T])\cap
C^0(\mathbb R^n\times[0,T])$ solve the heat equation and satisfy
$|u_i(x,t)|\leq Ae^{B|x|^2}$.  If $u_1(\cdot,0)=u_2(\cdot,0)$, then
$u_1=u_2$ on $\mathbb R^n\times[0,T]$.
\end{theorem}

\begin{proof}
\sketch
\uses{thm:euclidean-heat-wmp-exponential-quadratic-growth}
Apply Theorem~\ref{thm:euclidean-heat-wmp-exponential-quadratic-growth} to
$u_1-u_2$ and to $u_2-u_1$.
\end{proof}

\begin{theorem}[Nonuniqueness on $\mathbb R$ when bounded but not uniformly]
\label{thm:heat-nonuniqueness-slicewise-bounded}
\source{chow-et-al-ricci-flow-part-ii:page-0170}
\dcref{ch12:12.4}
\group{Weak Scalar Maximum Principles}

There is a nonzero solution $u:\mathbb R\times[0,\infty)\to\mathbb R$ of the
heat equation with $u(\cdot,0)=0$ such that, for every $\varepsilon>0$, a
constant $C(\varepsilon)<\infty$ satisfies
\[
|u(x,t)|\leq C(\varepsilon)e^{\varepsilon/t}
\qquad(x\in\mathbb R, t>0).
\]
Thus every positive-time slice is bounded, although no bound is uniform as
$t\downarrow0$.
\end{theorem}

\subsection{Constructing barrier functions}

\begin{lemma}[Distance-like function with first and second derivative control, II]
\label{lem:ricci-flow-distance-like-function}
\source{chow-et-al-ricci-flow-part-ii:page-0170}
\dcref{ch12:12.5}
\group{Barrier Functions}

Suppose $g(t)$, $0\leq t\leq T$, is a complete Ricci flow on a noncompact
$n$-manifold with $|\Rm(g(t))|\leq k_0$.  For every $O\in M$ there is a smooth
$h:M\to\mathbb R$ and $C=C(n,k_0,T)>0$ such that
\[
C^{-1}(d_{g(t)}(O,x)+1)\leq h(x)\leq C(d_{g(t)}(O,x)+1),
\]
\[
|\nabla h|_{g(t)}\leq C,\qquad \Hess_{g(t)}h\leq Cg(t)
\]
on $M\times[0,T]$.
\end{lemma}

\begin{proof}
\sketch
\uses{lem:bounded-curvature-distance-like-function}
\source{chow-et-al-ricci-flow-part-ii:page-0170,chow-et-al-ricci-flow-part-ii:page-0171}
Take the function supplied by
Lemma~\ref{lem:bounded-curvature-distance-like-function} for $g(0)$.  Bounded
Ricci curvature makes $g(t)$ and $g(0)$ uniformly equivalent, giving the
distance and gradient bounds.  Moreover
\[
\partial_t(\Hess_{g(t)}h)_{ij}
=(\nabla_iR^k{}_j+\nabla_jR^k{}_i-\nabla^kR_{ij})\nabla_kh.
\]
Shi's derivative estimate gives $|\nabla\Ric|\leq C t^{-1/2}$, while the
gradient is uniformly bounded.  Integrating this inequality from $0$ to $t$
and using uniform equivalence of the metrics yields the Hessian bound.
\end{proof}

\begin{remark}[Two-sided Hessian control in nonpositive curvature]
\label{rem:distance-like-hessian-lower-bound}
\source{chow-et-al-ricci-flow-part-ii:page-0170}
\uses{lem:ricci-flow-distance-like-function}
\dcref{ch12:12.6}
\group{Barrier Functions}


If $-b^2\leq\sec(g(t))\leq0$, Hessian comparison also gives a lower bound for
$\Hess h$.
\end{remark}

\begin{lemma}[Existence of a barrier supersolution, I]
\label{lem:general-metric-barrier-supersolution}
\source{chow-et-al-ricci-flow-part-ii:page-0172}
\dcref{ch12:12.7}
\group{Barrier Functions}

Let $\widetilde g$ be complete with $|\widetilde{\Rm}|\leq k_0$, let
$g(t)\geq\delta\widetilde g$, and let $|X(t)|_{\widetilde g}\leq b_1$.  For
every $\alpha,\beta>0$ there are a positive smooth $\phi$ and $\Lambda>0$ such
that
\[
(\partial_t-g^{ij}\widetilde\nabla_i\widetilde\nabla_j
-\widetilde\nabla_{X(t)})\phi\geq\beta\phi,
\]
\[
e^{\alpha\widetilde d(O,p)}\leq\phi(p,t)
 \leq e^{\Lambda(\widetilde d(O,p)+1)}.
\]
\end{lemma}

\begin{proof}
\sketch
\uses{lem:bounded-curvature-distance-like-function}
Let $h$ be given by Lemma~\ref{lem:bounded-curvature-distance-like-function}
and put $\phi=e^{Bt+\gamma h}$.  Since $g^{-1}\leq\delta^{-1}\widetilde
g^{-1}$,
\[
\phi^{-1}(\partial_t-g^{ij}\widetilde\nabla_i\widetilde\nabla_j
-\widetilde\nabla_X)\phi
\geq B-\delta^{-1}(\gamma nC+\gamma^2C^2)-b_1\gamma C.
\]
Choose $\gamma$ large enough for the required lower growth, then choose $B$
so that the right side is at least $\beta$.  The two-sided growth of $h$
gives the upper estimate.
\end{proof}

\begin{lemma}[Existence of a barrier supersolution, II]
\label{lem:ricci-flow-barrier-supersolution}
\source{chow-et-al-ricci-flow-part-ii:page-0173}
\dcref{ch12:12.9}
\group{Barrier Functions}

Let $g(t)$ be a complete Ricci flow with bounded curvature and let
$|X(t)|_{g(t)}\leq b_1$.  For every $a,A>0$ there are a positive smooth
$\phi$ and $b<\infty$ such that
\[
(\partial_t-\Delta-\nabla_{X(t)})\phi\geq A\phi,
\qquad
e^{a d_{g(t)}(O,x)}\leq\phi(x,t)\leq
e^{b(d_{g(t)}(O,x)+1)}.
\]
\end{lemma}

\begin{proof}
\sketch
\uses{lem:ricci-flow-distance-like-function}
For the function $h$ from Lemma~\ref{lem:ricci-flow-distance-like-function},
put $\phi=e^{Bt+\alpha h}$.  Choose $\alpha$ for the lower growth and then
$B$ so that
\[
B-\alpha\Delta h-\alpha^2|\nabla h|^2-
\alpha\langle X,\nabla h\rangle\geq A.
\]
The upper growth follows from the upper bound for $h$.
\end{proof}

\subsection{Scalar parabolic equations on complete manifolds}

\begin{theorem}[Weak maximum principle]
\label{thm:complete-manifold-scalar-wmp-barrier}
\source{chow-et-al-ricci-flow-part-ii:page-0173,chow-et-al-ricci-flow-part-ii:page-0174}
\dcref{ch12:12.10}
\group{Weak Scalar Maximum Principles}

Under the metric assumptions of Lemma~\ref{lem:general-metric-barrier-supersolution},
suppose $u\in C^{2,1}$, $u(\cdot,0)\leq0$, and
\[
(\partial_t-g^{ij}\widetilde\nabla_i\widetilde\nabla_j
-\widetilde\nabla_X)u\leq \frac A{\sqrt t}u
\quad\hbox{wherever }u\geq0.
\]
If $u(x,t)\leq e^{\widetilde b(\widetilde d(O,x)+1)}$, then $u\leq0$.
\end{theorem}

\begin{proof}
\sketch
\uses{lem:general-metric-barrier-supersolution}
Set $\widetilde u=e^{-2A\sqrt t}u$.  Then the displayed operator applied to
$\widetilde u$ is nonpositive wherever $\widetilde u\geq0$.  Let $\phi$ be
the barrier from Lemma~\ref{lem:general-metric-barrier-supersolution} with
growth exponent $2\widetilde b$ and differential lower bound $\phi$.  If
$\widetilde u-\varepsilon\phi$ first reached zero, its time derivative would
be nonnegative and its spatial Hessian nonpositive there, whereas the
differential inequality would make the parabolic operator strictly negative.
Thus $\widetilde u\leq\varepsilon\phi$ for every $\varepsilon>0$, and letting
$\varepsilon\downarrow0$ proves the result.
\end{proof}

\begin{example}[Integrable zeroth-order coefficient]
\label{ex:wmp-integrable-time-coefficient}
\source{chow-et-al-ricci-flow-part-ii:page-0175}
\dcref{ch12:12.11}
\group{Weak Scalar Maximum Principles}

The conclusion of Theorem~\ref{thm:complete-manifold-scalar-wmp-barrier}
remains true with $A/\sqrt t$ replaced by any integrable function
$\alpha:(0,T]\to\mathbb R$.
\end{example}

\begin{proof}
\sketch
\uses{thm:complete-manifold-scalar-wmp-barrier}
Replace $u$ by
$\exp(-\int_0^t\alpha(s)\,ds)u$.  The new function has no positive
zeroth-order term and the same spatial growth, so the barrier argument in
Theorem~\ref{thm:complete-manifold-scalar-wmp-barrier} applies verbatim.
\end{proof}

\begin{proposition}[Uniqueness for Ricci--DeTurck flow]
\label{prop:ricci-deturck-uniqueness-controlled-class}
\source{chow-et-al-ricci-flow-part-ii:page-0175}
\dcref{ch12:12.12}
\group{Ricci--DeTurck Flow}

Let $(M,\widetilde g)$ be complete with bounded curvature.  A solution $g(t)$
of Ricci--DeTurck flow with prescribed initial metric is unique among solutions
satisfying, for one finite $A$,
\[
A^{-1}\widetilde g\leq g(t)\leq A\widetilde g,
\qquad
|\widetilde\nabla g(t)|+\sqrt t,|\widetilde\nabla^2g(t)|\leq A.
\]
\end{proposition}

\begin{proof}
\sketch
\uses{thm:complete-manifold-scalar-wmp-barrier,ex:wmp-integrable-time-coefficient}
For two such solutions $g,\bar g$, the standard difference computation gives
\[
(\partial_t-g^{j\ell}\widetilde\nabla_j\widetilde\nabla_\ell)
|g-\bar g|^2\leq(Ct^{-1/2}+C^2)|g-\bar g|^2.
\]
The coefficient is integrable and the difference vanishes initially.
The integrable-coefficient estimate applied to the nonnegative
function $|g-\bar g|^2$ gives $g=\bar g$.
\end{proof}

\begin{remark}[Integrable time singularities]
\label{rem:wmp-integrable-time-singularity}
\source{chow-et-al-ricci-flow-part-ii:page-0175}
\uses{ex:wmp-integrable-time-coefficient}
\dcref{ch12:12.13}
\group{Weak Scalar Maximum Principles}


The factor $t^{-1/2}$ may be replaced by any function integrable on $(0,T]$.
\end{remark}

\begin{theorem}[Comparison with the associated ODE]
\label{thm:noncompact-scalar-pde-ode-comparison}
\source{chow-et-al-ricci-flow-part-ii:page-0175,chow-et-al-ricci-flow-part-ii:page-0176}
\dcref{ch12:12.14}
\group{Weak Scalar Maximum Principles}

Let $g(t)$ be a Ricci flow.  Assume it is complete with bounded curvature, and
\[
Lu=u_t-\Delta u-\langle X,\nabla u\rangle-F(u,t),
\]
where $X$ is bounded and $F$ is continuous and locally Lipschitz in its first
variable.  Suppose $Lu\leq0$ and
$|u(x,t)|\leq e^{\widetilde b(d_{g(t)}(O,x)+1)}$.  If $U'=F(U,t)$ and
$u(\cdot,0)\leq U(0)$, then $u(x,t)\leq U(t)$ as long as $U$ exists.
\end{theorem}

\begin{proof}
\sketch
\uses{lem:ricci-flow-barrier-supersolution}
For $v=u-U$ and $\widetilde v=e^{-A_1t}v$, use the barrier $\phi$ from
Lemma~\ref{lem:ricci-flow-barrier-supersolution}.  If
$\widetilde v-\varepsilon\phi$ first vanished at a positive maximum, local
Lipschitz continuity gives $F(u,t)-F(U,t)\leq K_0v$ there.  Hence the
parabolic operator applied to the difference is at most
$(K_0-A_1)\varepsilon\phi$.  Taking $A_1>K_0$ contradicts the derivative
test.  Letting $\varepsilon\downarrow0$ yields $v\leq0$.
\end{proof}

\subsection{Weak solutions on complete manifolds}

\begin{definition}[Continuous subharmonic functions]
\label{def:continuous-subharmonic-function}
\source{chow-et-al-ricci-flow-part-ii:page-0178}
\dcref{ch12:12.15}
\group{Weak Scalar Maximum Principles}

A continuous function $u$ on an open set $\Omega$ is subharmonic if, for every
relatively compact domain $D\Subset\Omega$ and every harmonic $h$ on $D$ with
$u\leq h$ on $\partial D$, one has $u\leq h$ on $D$.  It is superharmonic if
$-u$ is subharmonic and harmonic if it is both.
\end{definition}

\begin{lemma}[Absolute value of a solution is a subsolution]
\label{lem:absolute-value-harmonic-subsolution}
\source{chow-et-al-ricci-flow-part-ii:page-0178}
\dcref{ch12:12.16}
\group{Weak Scalar Maximum Principles}

If $u$ is continuous and harmonic, then $u_+=\max\{u,0\}$,
$u_-=\max\{-u,0\}$, and $|u|$ are nonnegative continuous subharmonic
functions.
\end{lemma}

\begin{proof}
\sketch
\uses{def:continuous-subharmonic-function}
The maximum of two continuous subharmonic functions is subharmonic: if a
harmonic $h$ dominates both on the boundary of a relatively compact domain,
the comparison definition makes it dominate both in the interior.  Apply this
to $(u,0)$ and $(-u,0)$; their sum is $|u|$.
\end{proof}

\begin{remark}[Minimum of superharmonic functions]
\label{rem:minimum-superharmonic-functions}
\source{chow-et-al-ricci-flow-part-ii:page-0178}
\uses{lem:absolute-value-harmonic-subsolution}
\dcref{ch12:12.17}
\group{Weak Scalar Maximum Principles}


The minimum of two continuous superharmonic functions is superharmonic.  The
same lattice property holds for subsolutions and supersolutions of the heat
equation.
\end{remark}

\begin{theorem}[WMP for weak subsolutions on noncompact manifolds]
\label{thm:static-weak-subsolution-wmp}
\source{chow-et-al-ricci-flow-part-ii:page-0178,chow-et-al-ricci-flow-part-ii:page-0179}
\dcref{ch12:12.18}
\group{Weighted Energy Maximum Principles}

Let $(M,g)$ be complete.  If $u$ is a weak subsolution of $u_t-\Delta u=0$,
$u(\cdot,0)\leq0$, and for some $\alpha>0$
\[
\int_0^T\!\int_M e^{-\alpha d(x,O)^2}u_+(x,t)^2\,d\mu\,dt<\infty,
\]
then $u\leq0$ on $M\times[0,T]$.
\end{theorem}

\begin{proof}
\sketch
\uses{thm:evolving-metric-weak-subsolution-wmp}
Apply Theorem~\ref{thm:evolving-metric-weak-subsolution-wmp} to the constant
metric family, with $g_*=g$ and $\mathcal H=0$.
\end{proof}

\begin{definition}[Ricci curvature with a negative quadratic lower bound]
\label{def:negative-quadratic-ricci-lower-bound}
\source{chow-et-al-ricci-flow-part-ii:page-0179}
\dcref{ch12:12.19}
\group{Weighted Energy Maximum Principles}

A noncompact Riemannian manifold has Ricci curvature with a negative
quadratic lower bound if, for some $O\in M$ and $C<\infty$,
\[
\Ric(x)\geq-C(d(x,O)^2+1)g_x.
\]
\end{definition}

\begin{corollary}[WMP for bounded weak subsolutions]
\label{cor:bounded-weak-subsolution-wmp}
\source{chow-et-al-ricci-flow-part-ii:page-0180}
\dcref{ch12:12.20}
\group{Weighted Energy Maximum Principles}

On a complete noncompact manifold with a negative quadratic Ricci lower
bound, every bounded weak heat subsolution with nonpositive initial data is
nonpositive.  Consequently bounded weak heat solutions are unique.
\end{corollary}

\begin{proof}
\sketch
\uses{def:negative-quadratic-ricci-lower-bound,thm:static-weak-subsolution-wmp}
Bishop--Gromov comparison gives $\Vol B(O,r)\leq ae^{br^2}$.  If $|u|\leq C_1$,
decomposition into unit
annuli gives
\[
\int_0^T\!\int_M e^{-\alpha d^2}u_+^2
\leq aC_1^2T\sum_{r\geq1}e^{br^2-\alpha(r-1)^2}<\infty
\]
for $\alpha>b$.  Theorem~\ref{thm:static-weak-subsolution-wmp} applies.
Uniqueness follows by applying it to both signs of the difference.
\end{proof}

\begin{remark}[Li--Yau uniqueness]
\label{rem:li-yau-below-bounded-uniqueness}
\source{chow-et-al-ricci-flow-part-ii:page-0180}
\uses{def:negative-quadratic-ricci-lower-bound}
\dcref{ch12:12.21}
\group{Weighted Energy Maximum Principles}


Under the same Ricci lower bound, Li--Yau gradient estimates give uniqueness
even for heat solutions bounded only from below.
\end{remark}

\begin{theorem}[WMP on noncompact manifolds with evolving metrics]
\label{thm:evolving-metric-weak-subsolution-wmp}
\source{chow-et-al-ricci-flow-part-ii:page-0180,chow-et-al-ricci-flow-part-ii:page-0181}
\dcref{ch12:12.22}
\group{Weighted Energy Maximum Principles}

Let $g_t=-2h$ be a family of complete metrics with $g(t)\geq g_*$ for a
complete metric $g_*$, and put $\mathcal H=\operatorname{tr}_{g(t)}h$.  Assume
$\mathcal H_*(t)=\inf_M\mathcal H(\cdot,t)$ is integrable.  If $u$ is a weak
subsolution of $u_t-\Delta_{g(t)}u=0$, $u(\cdot,0)\leq0$, and
\[
\int_0^T\!\int_Me^{-bd_*^2(O,x)}u_+^2\,d\mu_t\,dt<\infty,
\]
then $u\leq0$ on $M\times[0,T]$.
\end{theorem}

\begin{proof}
\sketch
\source{chow-et-al-ricci-flow-part-ii:page-0181,chow-et-al-ricci-flow-part-ii:page-0182,chow-et-al-ricci-flow-part-ii:page-0183}
It suffices to treat $T\leq(8b)^{-1}$.  For $T_1\leq T$ set
\[
\eta=-\frac{d_*^2(O,x)}{4(2T_1-t)},\qquad
\widetilde{\mathcal H}(t)=
\exp\!\left(\int_0^t\mathcal H_*(s)\,ds\right).
\]
Since $g(t)\geq g_*$,
$\eta_t+|\nabla\eta|_{g(t)}^2\leq0$.  Choose a radial cutoff $\varphi_r$ that
is one on $B_*(O,r/2)$, vanishes outside $B_*(O,r)$, and satisfies
$|\nabla\varphi_r|^2\leq9r^{-2}$.  Testing the weak inequality with
$\widetilde{\mathcal H}\varphi_r^2e^\eta u_+$, integrating by parts, and
completing squares gives
\[
\frac12\int_M\widetilde{\mathcal H}(T_1)\varphi_r^2e^\eta
u_+(T_1)^2\,d\mu_{T_1}
\leq2\int_0^{T_1}\!\int_M
\widetilde{\mathcal H}|\nabla\varphi_r|^2e^\eta u_+^2\,d\mu_t\,dt.
\]
The assumed weighted integrability and dominated convergence make the right
side tend to zero as $r\to\infty$.  Thus $u_+(\cdot,T_1)=0$.  Since $T_1$ is
arbitrary, and successive short time intervals cover $[0,T]$, the result
follows.
\end{proof}

\begin{remark}[Positive part suffices]
\label{rem:evolving-wmp-positive-part}
\source{chow-et-al-ricci-flow-part-ii:page-0183}
\uses{thm:evolving-metric-weak-subsolution-wmp}
\dcref{ch12:12.23}
\group{Weighted Energy Maximum Principles}


The conclusion remains valid if only $u_+$ is assumed to be a weak
subsolution.
\end{remark}

\begin{theorem}[$L^1$ vanishing]
\label{thm:l1-heat-subsolution-vanishing}
\source{chow-et-al-ricci-flow-part-ii:page-0183}
\dcref{ch12:12.24}
\group{Heat Equation Uniqueness}

Let $(M,g)$ be complete, noncompact, and have a negative quadratic Ricci lower
bound.  If $u\geq0$ is a heat subsolution,
$u(\cdot,t)\in L^1(M)$ for every $t>0$, and
$\lim_{t\downarrow0}\int_Mu(x,t)\,d\mu=0$, then $u=0$ on
$M\times(0,\infty)$.
\end{theorem}

\begin{corollary}[$L^1$ uniqueness]
\label{cor:l1-heat-uniqueness}
\source{chow-et-al-ricci-flow-part-ii:page-0183}
\uses{lem:absolute-value-harmonic-subsolution, thm:l1-heat-subsolution-vanishing}
\dcref{ch12:12.25}
\group{Heat Equation Uniqueness}


Under the assumptions of Theorem~\ref{thm:l1-heat-subsolution-vanishing}, two
$C^{2,1}$ heat solutions in $L^1(M)$ with equal initial data are equal for all
positive times.
\end{corollary}

\begin{proof}
\sketch
The absolute value of their difference is a nonnegative heat subsolution;
apply Theorem~\ref{thm:l1-heat-subsolution-vanishing}.
\end{proof}

\begin{theorem}[$L^p$ uniqueness for $p>1$]
\label{thm:lp-heat-subsolution-vanishing}
\source{chow-et-al-ricci-flow-part-ii:page-0183}
\dcref{ch12:12.26}
\group{Heat Equation Uniqueness}

Let $(M,g)$ be any complete noncompact manifold and $1<p<\infty$.  If
$u\geq0$ is a heat subsolution,
$\int_0^{T'}\int_Mu^p<\infty$ for every $T'<T$, and
$\lim_{t\downarrow0}\int_Mu(x,t)^p\,d\mu=0$, then
$u=0$ on $M\times[0,T)$.
\end{theorem}

\begin{proof}
\sketch
\source{chow-et-al-ricci-flow-part-ii:page-0184}
Let $\varphi_r$ be supported in $B(O,2r)$, equal to one on $B(O,r)$, and
$|\nabla\varphi_r|^2\leq3r^{-2}$.  Multiplying the subsolution inequality by
$p\varphi_r^2u^{p-1}$ and integrating by parts and in time gives
\[
\int_M\varphi_r^2u(T')^p+
\frac{2(p-1)}p\int_0^{T'}\!\int_M
\varphi_r^2|\nabla u^{p/2}|^2
\leq\frac{6p}{(p-1)r^2}\int_0^{T'}\!\int_Mu^p.
\]
Letting $r\to\infty$ makes the right side vanish and proves $u(T')=0$.
\end{proof}

\section{Mollifying distance functions on Riemannian manifolds}

Let $\eta$ be the standard radial unit-mass mollifier on $\mathbb R^n$.  If
$f$ is locally integrable and $\varepsilon$ is below the conjugate-radius
scale, define
\[
f_\varepsilon(p)=\varepsilon^{-n}\int_{B(\varepsilon)}
\eta(v/\varepsilon)f(\exp_p\iota_p(v))\,dv,
\]
where $\iota_p:\mathbb R^n\to T_pM$ is an isometry.

\begin{lemma}[Smoothness of intrinsic mollification]
\label{lem:intrinsic-mollification-smooth}
\source{chow-et-al-ricci-flow-part-ii:page-0187}
\dcref{ch12:12.28}
\group{Mollified Distance Functions}

Let $(M,g)$ be complete with sectional curvature bounded above by $k_0$ and
let $f$ be locally integrable on $B(O,r_0)$.  Put
\[
\widehat r_0(p)=\sup\{r\leq\pi/\sqrt{k_0}:B(p,r)\subset B(O,r_0)\}.
\]
For $p\in B(O,r_0-\varepsilon)$ and
$0<\varepsilon<\widehat r_0(p)/10$, intrinsic mollification commutes with the
local exponential pullback, and both the pulled-back mollification and
$f_\varepsilon$ are smooth on their natural domains.
\end{lemma}

\begin{proof}
\sketch
\source{chow-et-al-ricci-flow-part-ii:page-0187,chow-et-al-ricci-flow-part-ii:page-0188}
On the exponential pullback ball, local isometry and radiality of $\eta$
identify the two convolution formulas.  After changing variables, the pulled
back formula is an integral of $f$ against
\[
\eta\!\left(\varepsilon^{-1}(\widehat{\exp}_q\widehat{\iota}_q)^{-1}(y)\right)
 J_{(\widehat{\exp}_q\widehat{\iota}_q)^{-1}}(y).
\]
The inverse exponential and its Jacobian are smooth in $(q,y)$; the kernel
has compact support, so differentiation under the integral proves smoothness.
Composing with the inverse of $\exp_p$ on an injectivity neighborhood proves
smoothness on $M$.
\end{proof}

\begin{lemma}[Mollifiers and derivatives]
\label{lem:intrinsic-mollifier-derivative-convergence}
\source{chow-et-al-ricci-flow-part-ii:page-0188,chow-et-al-ricci-flow-part-ii:page-0189}
\dcref{ch12:12.29}
\group{Mollified Distance Functions}

Under the preceding assumptions, in normal coordinates at $p$:
\begin{enumerate}
\item if $f$ is continuous at $p$, then $f_\varepsilon(p)\to f(p)$, uniformly
on compact sets where $f$ is continuous;
\item if $f$ has a first Aleksandrov derivative at $p$, then
$\partial_if_\varepsilon(p)\to\partial_if(p)$;
\item if $f$ has a second Aleksandrov derivative at $p$, then
$\partial_i\partial_jf_\varepsilon(p)\to\partial_i\partial_jf(p)$;
\item if $f\in L^p_{\mathrm{loc}}$, $1\leq p<\infty$, then
$f_\varepsilon\to f$ in $L^p_{\mathrm{loc}}$.
\end{enumerate}
\end{lemma}

\begin{proof}
\sketch
\uses{lem:intrinsic-mollification-smooth}
\source{chow-et-al-ricci-flow-part-ii:page-0189,chow-et-al-ricci-flow-part-ii:page-0190,chow-et-al-ricci-flow-part-ii:page-0191,chow-et-al-ricci-flow-part-ii:page-0192,chow-et-al-ricci-flow-part-ii:page-0193}
The first assertion follows by bounding the convolution error by the modulus
of continuity on an $\varepsilon$-ball.  For the derivative assertions, write
the first- or second-order Aleksandrov expansion of $f$ at $p$.  Differentiate
the smooth kernel from Lemma~\ref{lem:intrinsic-mollification-smooth}.  Its
unit mass kills the constant term, its first moments reproduce the linear
term, and its second moments reproduce the quadratic term; after the scaling
$v=\varepsilon w$, the Aleksandrov remainder is bounded by its defining
$o(|v|)$ or $o(|v|^2)$ estimate times an integrable fixed kernel.  Dominated
convergence therefore gives (ii) and (iii).  Finally, on each compact set the
intrinsic kernels form a uniformly bounded approximate identity in finitely
many normal-coordinate charts.  The usual density argument, first for
continuous functions and then for $L^p$, proves (iv).
\end{proof}

\begin{lemma}[Distance-like function with first and second derivative control, I]
\label{lem:bounded-curvature-distance-like-function}
\source{chow-et-al-ricci-flow-part-ii:page-0193,chow-et-al-ricci-flow-part-ii:page-0194}
\dcref{ch12:12.30}
\group{Mollified Distance Functions}

Let $(M,g)$ be complete and noncompact with $|\Rm|\leq k_0$.  For every
$O\in M$ there is a smooth $h:M\to\mathbb R$ and $C=C(n,k_0)$ such that
\[
C^{-1}(d(O,x)+1)\leq h(x)\leq C(d(O,x)+1),
\qquad |\nabla h|\leq C,qquad \Hess h\leq Cg.
\]
\end{lemma}

\begin{proof}
\sketch
\uses{lem:intrinsic-mollification-smooth}
\source{chow-et-al-ricci-flow-part-ii:page-0194,chow-et-al-ricci-flow-part-ii:page-0195,chow-et-al-ricci-flow-part-ii:page-0196,chow-et-al-ricci-flow-part-ii:page-0197}
Mollify $f=d(O,\cdot)+2$ at one fixed sufficiently small scale
$\varepsilon$.  The triangle inequality and unit mass of the kernel give
$d(O,x)+1\leq f_\varepsilon(x)\leq d(O,x)+3$.  To differentiate, identify
nearby tangent spaces by parallel translation.  A directional derivative is
the average of $\langle\nabla f,J\rangle$, where $J$ is the Jacobi field with
$J(0)=V$ and $J'(0)=0$.  Rauch comparison bounds $J$, proving the gradient
bound.  A second differentiation gives
\[
\Hess f_\varepsilon(V,V)=\varepsilon^{-n}\int
\bigl(\Hess f(J,J)+\langle\nabla f,\nabla_JJ\rangle\bigr)
\eta(|v|/\varepsilon)\,dv.
\]
Hessian comparison bounds the first term from above.  For the geodesic
variation defining $J$, curvature symmetry and $\nabla_TJ=\nabla_JT$ give
\[
\langle T,\nabla_JJ\rangle
=\int\langle T,R(T,J)J\rangle-\int|\nabla_TJ|^2.
\]
Bounded curvature and the Jacobi equation uniformly bound this expression.
Thus the Hessian integral has a uniform upper bound, completing the proof.
\end{proof}
\section{Weak maximum principles for parabolic systems}

Let $\mathcal W\to M$ be a metric vector bundle with compatible
time-dependent connections, and let
$\mathcal V=\mathcal W\otimes_{\mathrm S}\mathcal W$.  We consider symmetric
bilinear-form supersolutions
\[
u_t\geq\Delta u+\nabla_{X(t)}u+\beta(u,t).
\]

\begin{proposition}[WMP for systems: preserving nonnegativity]
\label{prop:noncompact-system-preserves-nonnegativity}
\source{chow-et-al-ricci-flow-part-ii:page-0197,chow-et-al-ricci-flow-part-ii:page-0198}
\dcref{ch12:12.31}
\group{Weak System Maximum Principles}

Suppose $(M,g(t))$ is complete and noncompact with uniformly bounded
curvature.  Assume $\beta$ is continuous in time, locally Lipschitz in $u$,
and satisfies the null-eigenvector condition
\[
\omega\geq0,\quad \omega(V,V)=0
\quad\Longrightarrow\quad \beta(\omega,t)(V,V)\geq0.
\]
If $u(0)\geq0$ and
$|u(x,t)|\leq e^{\widehat a(d_{g(t)}(O,x)+1)}$, then $u(t)\geq0$ for all
$t\in[0,T]$.
\end{proposition}

\begin{proof}
\sketch
\uses{lem:ricci-flow-barrier-supersolution}
\source{chow-et-al-ricci-flow-part-ii:page-0198,chow-et-al-ricci-flow-part-ii:page-0199,chow-et-al-ricci-flow-part-ii:page-0200}
Let $\widetilde h$ be the identity positive form on $\mathcal W$ and let
$\phi$ be the barrier in Lemma~\ref{lem:ricci-flow-barrier-supersolution},
with growth exponent larger than that of $u$.  For $\varepsilon>0$ put
$\widetilde u=u+\varepsilon\phi\widetilde h$.  If positivity first failed at
$(x_0,t_0)$ in the unit direction $V_0$, the growth of $\phi$ would keep
$x_0$ in a fixed compact set.  Parallel extend $V_0$ so that its first and
second covariant derivatives vanish at $x_0$.  The derivative test and the
equation then give
\[
0\geq \beta(\widetilde u,t_0)(V_0,V_0)
+\varepsilon(A-K_0-C_2)\phi(x_0,t_0),
\]
where $K_0$ is a Lipschitz constant for $\beta$ on that compact set and
$|\langle X,\nabla\phi\rangle|\leq C_2\phi$.  The first term is nonnegative
by the null-eigenvector condition.  Choosing $A>K_0+C_2$ is contradictory.
Thus $\widetilde u>0$; let $\varepsilon\downarrow0$.
\end{proof}

\begin{remark}[Closed-manifold case]
\label{rem:system-wmp-closed-manifolds}
\source{chow-et-al-ricci-flow-part-ii:page-0198}
\uses{prop:noncompact-system-preserves-nonnegativity}
\dcref{ch12:12.32}
\group{Weak System Maximum Principles}


The same proof without a barrier at infinity gives the corresponding result
on closed manifolds.
\end{remark}

\begin{theorem}[Maximum principle for symmetric $2$-tensors]
\label{thm:noncompact-two-tensor-maximum-principle}
\source{chow-et-al-ricci-flow-part-ii:page-0198}
\dcref{ch12:12.33}
\group{Weak System Maximum Principles}

Let $(M,g(t))$ be complete, noncompact, and of bounded curvature.  Suppose a
symmetric $2$-tensor $\alpha$ satisfies
\[
\alpha_t\geq\Delta\alpha+\nabla_X\alpha+\beta(\alpha,t),
\qquad
|\alpha(x,t)|\leq e^{\widehat a(d_{g(t)}(O,x)+1)},
\]
where $\beta$ is continuous in time, locally Lipschitz in $\alpha$, and
satisfies the null-eigenvector condition.  If $\alpha(0)\geq0$, then
$\alpha(t)\geq0$ on $[0,T]$.
\end{theorem}

\begin{proof}
\sketch
\uses{prop:noncompact-system-preserves-nonnegativity}
Apply Proposition~\ref{prop:noncompact-system-preserves-nonnegativity} to
$\mathcal W=T^*M$ and $u=\alpha$.
\end{proof}

\subsection{Weinberger--Hamilton maximum principles}

\begin{theorem}[Weinberger--Hamilton maximum principle for systems]
\label{thm:weinberger-hamilton-noncompact-systems}
\source{chow-et-al-ricci-flow-part-ii:page-0200}
\dcref{ch12:12.34}
\group{Weak System Maximum Principles}

Let $\mathcal K\subset\mathcal V$ be closed, invariant under parallel
translation, and fiberwise convex.  Suppose $F(u,t)$ is continuous and locally
Lipschitz in $u$, and every solution of $U'=F_x(U,t)$ starting in
$\mathcal K_x$ remains there.  If
\[
u_t=\Delta u+\nabla_{X(t)}u+F(u,t),qquad
d(u(x,t),\mathcal K_x)\leq e^{\widehat a(d_{g(t)}(O,x)+1)},
\]
and $u(x,0)\in\mathcal K_x$, then $u(x,t)\in\mathcal K_x$ for all $x,t$.
\end{theorem}

\begin{proof}
\sketch
\uses{thm:time-dependent-convex-set-maximum-principle}
Take the constant family $\mathcal K(t)=\mathcal K$ in
Theorem~\ref{thm:time-dependent-convex-set-maximum-principle}.
\end{proof}

\begin{theorem}[Maximum principle for systems with time-dependent sets]
\label{thm:time-dependent-convex-set-maximum-principle}
\source{chow-et-al-ricci-flow-part-ii:page-0200,chow-et-al-ricci-flow-part-ii:page-0201}
\dcref{ch12:12.35}
\group{Weak System Maximum Principles}

For each $t$, let $\mathcal K(t)\subset\mathcal V$ be parallel-translation
invariant with closed convex fibers, and assume its space-time track is
closed.  Let $F$ be continuous and locally Lipschitz in $u$, and suppose every
solution of $U'=F_x(U,t)$ starting in $\mathcal K(t_0)_x$ remains in
$\mathcal K(t)_x$.  If a solution
$u_t=\Delta u+\nabla_Xu+F(u,t)$ satisfies
\[
d(u(x,t),\mathcal K(t)_x)\leq
e^{\widehat a(d_{g(t)}(O,x)+1)}
\]
and begins in $\mathcal K(0)$, then it remains in $\mathcal K(t)$.
\end{theorem}

\begin{proof}
\sketch
\uses{lem:ricci-flow-barrier-supersolution}
\source{chow-et-al-ricci-flow-part-ii:page-0201,chow-et-al-ricci-flow-part-ii:page-0202,chow-et-al-ricci-flow-part-ii:page-0203,chow-et-al-ricci-flow-part-ii:page-0204}
Write $\rho_t(v)=d(v,\mathcal K(t)_x)$ and
\[
s(t)=\sup_x\{\rho_t(u(x,t))-\varepsilon\phi(x,t)\},
\]
where $\phi$ is the barrier of
Lemma~\ref{lem:ricci-flow-barrier-supersolution}.  Its growth makes every
positive supremum occur in a fixed compact set.  At such a maximizing point,
represent distance to the closed convex fiber by a supporting functional
$\ell$ at the nearest point $v$.  Advance $v$ for a short time by the ODE;
ODE invariance keeps the advanced point in the next convex set.  Comparing
the support functional at the two times, using the PDE, and parallel
extending $(v,\ell)$ in space gives
\[
\frac{d^+s}{dt}\leq
\ell(\Delta u)+\ell(\nabla_Xu)+
\ell(F(u,t)-F(v,t))-\varepsilon\phi_t.
\]
At the spatial maximum, the Laplacian and drift terms are bounded by their
corresponding barrier terms.  Local Lipschitz continuity of $F$ bounds the
remaining difference by $C\rho_t(u)$, while
$\phi_t-\Delta\phi-\nabla_X\phi\geq A\phi$.  Choosing $A>C$ yields
$d^+s/dt\leq Cs$.  Since $s(0)\leq0$, the upper-Dini-derivative comparison
implies $s\leq0$.  Let $\varepsilon\downarrow0$.
\end{proof}

\begin{remark}[Applying the barrier to distance]
\label{rem:barrier-applied-to-distance-from-convex-set}
\source{chow-et-al-ricci-flow-part-ii:page-0201}
\uses{thm:time-dependent-convex-set-maximum-principle}
\dcref{ch12:12.36}
\group{Weak System Maximum Principles}


The essential device is to apply the scalar barrier to
$d(u(x,t),\mathcal K(t)_x)$ rather than to the bundle section $u$ itself.
\end{remark}

\begin{proof}
\sketch
\uses{lem:ricci-flow-barrier-supersolution}
\source{chow-et-al-ricci-flow-part-ii:page-0204,chow-et-al-ricci-flow-part-ii:page-0205,chow-et-al-ricci-flow-part-ii:page-0206}
At a point $(x_0,t_0)$ maximizing
$\rho_{t_0}(u)-\varepsilon\phi$, compare with time $t_0-h$.  Convexity of
$\rho_{t_0-h}$ and the second derivative test imply that the contribution of
$-h\Delta u$ is at least $-h\varepsilon\Delta\phi+o(h)$.  If $v_h$ is the
nearest point in $\mathcal K(t_0-h)_{x_0}$, the ODE solution from $v_h$ lies
in $\mathcal K(t_0)_{x_0}$ and equals $v_h+hF(v_h)+o(h)$.  The Lipschitz
property of $F$ therefore gives
\[
\rho_{t_0}(u)-\rho_{t_0-h}(u-hF(u))
\leq hC\rho_{t_0}(u)+o(h).
\]
After the barrier terms are restored, $d^+s/dt\leq Cs$; since $s(0)=0$, one
gets $s\leq0$, and then lets $\varepsilon\downarrow0$.
\end{proof}

\begin{theorem}[Maximum principle for systems with avoidance sets]
\label{thm:noncompact-system-avoidance-set-maximum-principle}
\source{chow-et-al-ricci-flow-part-ii:page-0206,chow-et-al-ricci-flow-part-ii:page-0207}
\dcref{ch12:12.38}
\group{Weak System Maximum Principles}

Under the hypotheses of
Theorem~\ref{thm:time-dependent-convex-set-maximum-principle}, let
$\mathcal A(t)\subset\mathcal K(t)$.  Suppose each ODE solution starting in
$\mathcal K(t_0)_x$ either stays in $\mathcal K(t)_x$, or remains in
$\mathcal K(t)_x\setminus\mathcal A(t)_x$ until a first time when it reaches
$\mathcal A(t)_x$.  A PDE solution with the same growth bound that starts in
$\mathcal K(0)$ and avoids $\mathcal A(t)$ remains in $\mathcal K(t)$.
\end{theorem}

\begin{proof}
\sketch
\uses{thm:time-dependent-convex-set-maximum-principle}
Run the proof of
Theorem~\ref{thm:time-dependent-convex-set-maximum-principle} up to a
hypothetical first escape.  The short ODE comparison used there cannot meet
$\mathcal A$, because the PDE solution and the compact maximizing sequence
avoid $\mathcal A$ and the space-time track is closed.  It therefore stays in
$\mathcal K$, so the same supporting-functional estimate yields the same
contradiction.
\end{proof}

\begin{remark}[Time-dependent bundle metrics]
\label{rem:maximum-principles-time-dependent-bundle-metric}
\source{chow-et-al-ricci-flow-part-ii:page-0207}
\uses{thm:noncompact-two-tensor-maximum-principle, thm:weinberger-hamilton-noncompact-systems, thm:time-dependent-convex-set-maximum-principle, thm:noncompact-system-avoidance-set-maximum-principle}
\dcref{ch12:12.39}
\group{Weak System Maximum Principles}


The system maximum principles remain valid for a smooth family of compatible
bundle metrics $h(t)$ satisfying $\partial_th\geq-c_1h$, after adding the
resulting zeroth-order terms to the barrier estimates.  This formulation can
replace Uhlenbeck's trick in curvature-operator applications.
\end{remark}

\section{Strong maximum principles for parabolic systems}

\subsection{Scalar heat-type equations}

For
\[
Lu=\Delta_{g(t)}u+\nabla_{X(t)}u+c(x,t)u-u_t
\]
on a space-time domain $D$, let $\gamma(D)$ be its upper boundary and let
$S(x_0,t_0)$ consist of points joined to $(x_0,t_0)$ by a curve in $D$ whose
time coordinate is nondecreasing toward $(x_0,t_0)$.

\begin{theorem}[Strong maximum principle for scalar equations]
\label{thm:manifold-parabolic-strong-maximum-principle}
\source{chow-et-al-ricci-flow-part-ii:page-0208}
\dcref{ch12:12.40}
\group{Strong Scalar Maximum Principles}

Assume $c\leq0$.  If $u\in C^{2,1}(D)\cap C^0(D\cup\gamma(D))$ satisfies
$Lu\geq0$ and attains a nonnegative maximum at
$(x_0,t_0)\in D\cup\gamma(D)$, then
\[
u(x,t)=u(x_0,t_0)qquad ((x,t)\in\overline{S(x_0,t_0)}).
\]
\end{theorem}

\begin{proof}
\sketch
\uses{thm:euclidean-parabolic-strong-maximum-principle}
Choose coordinate cylinders along any admissible curve to $(x_0,t_0)$.
Theorem~\ref{thm:euclidean-parabolic-strong-maximum-principle} makes the
maximum set backward-open in each cylinder; continuity makes it closed.
The connected chain of cylinders therefore lies in the maximum set.  Taking
closures proves the assertion.
\end{proof}

\begin{theorem}[Strong maximum principle for Euclidean scalar equations]
\label{thm:euclidean-parabolic-strong-maximum-principle}
\source{chow-et-al-ricci-flow-part-ii:page-0208}
\dcref{ch12:12.41}
\group{Strong Scalar Maximum Principles}

Let
\[
L_1u=a^{ij}u_{ij}+b^iu_i+cu-u_t
\]
on a domain $D_1\subset\mathbb R^n\times\mathbb R$, where the coefficients
are continuous, $(a^{ij})$ is positive definite, and $c\leq0$.  If a
$C^{2,1}$ subsolution $L_1u\geq0$ attains a nonnegative maximum at
$(x_0,t_0)\in D_1\cup\gamma(D_1)$, then it is constant on
$\overline{S(x_0,t_0)}$.
\end{theorem}

\begin{corollary}[Cylindrical strong maximum principle]
\label{cor:cylindrical-parabolic-strong-maximum-principle}
\source{chow-et-al-ricci-flow-part-ii:page-0209}
\uses{thm:euclidean-parabolic-strong-maximum-principle}
\dcref{ch12:12.42}
\group{Strong Scalar Maximum Principles}


If $\Omega\subset\mathbb R^n$ is connected and a subsolution as in
Theorem~\ref{thm:euclidean-parabolic-strong-maximum-principle} attains a
nonnegative maximum at $(x_0,t_0)\in\Omega\times(0,T]$, then it equals that
maximum on $\overline\Omega\times(0,t_0]$.
\end{corollary}

\begin{proof}
\sketch
Every point of $\Omega\times(0,t_0]$ belongs to the closure of the backward
reachable set of $(x_0,t_0)$; apply
Theorem~\ref{thm:euclidean-parabolic-strong-maximum-principle}.
\end{proof}

\begin{corollary}[Strong maximum principle for scalar curvature]
\label{cor:ricci-flow-scalar-curvature-strong-maximum-principle}
\source{chow-et-al-ricci-flow-part-ii:page-0209}
\uses{thm:manifold-parabolic-strong-maximum-principle}
\dcref{ch12:12.43}
\group{Ricci Flow Applications}


Let $g(t)$ be a Ricci flow on a connected manifold, not necessarily complete,
with $R\geq0$.  If $R(x_0,t_0)=0$ for some $t_0>0$, then
$\Ric=0$ on $M\times(0,t_0]$.  If also $\Rm\geq0$, then $\Rm=0$ there.
\end{corollary}

\begin{proof}
\sketch
Since $R_t=\Delta R+2|\Ric|^2$, the function $-R$ is a subsolution attaining
its nonnegative maximum zero.  Theorem~\ref{thm:manifold-parabolic-strong-maximum-principle}
makes $R$ vanish backward in time; its evolution equation then gives
$\Ric=0$.  A nonnegative curvature operator of trace zero is zero.
\end{proof}

For the Euclidean backward heat kernel $H(x_0,y,t_0-t)$, define the heat ball
\[
E_r(x_0,t_0)=\{(y,t):H(x_0,y,t_0-t)>r^{-n}\}.
\]

\begin{corollary}[Euclidean mean value property]
\label{cor:euclidean-heat-mean-value-property}
\source{chow-et-al-ricci-flow-part-ii:page-0210}
\dcref{ch12:12.44}
\group{Strong Scalar Maximum Principles}

If $u_t-\Delta_0u=0$ and $E_r(x_0,t_0)$ lies in its domain, then
\[
u(x_0,t_0)=\frac1{4r^n}\iint_{E_r(x_0,t_0)}
u(y,t)\frac{|x_0-y|^2}{(t_0-t)^2}\,dy\,dt.
\]
Moreover
\[
E_r=\left\{(y,t):|y-x_0|^2<2n(t_0-t)
\log\frac{r^2}{4\pi(t_0-t)}\right\}.
\]
\end{corollary}

\begin{remark}[Divergence-form mean value properties]
\label{rem:divergence-form-mean-value-property}
\source{chow-et-al-ricci-flow-part-ii:page-0210}
\uses{cor:euclidean-heat-mean-value-property}
\dcref{ch12:12.45}
\group{Strong Scalar Maximum Principles}


The heat-ball mean value property extends to uniformly parabolic operators in
divergence form.
\end{remark}

\begin{corollary}[Mean value strong maximum principle]
\label{cor:heat-mean-value-strong-maximum-principle}
\source{chow-et-al-ricci-flow-part-ii:page-0210}
\uses{cor:euclidean-heat-mean-value-property}
\dcref{ch12:12.46}
\group{Strong Scalar Maximum Principles}


If a heat solution on a connected cylinder attains its supremum at an
interior point $(x_1,t_1)$, it equals that supremum at all $(y,s)$ with
$s\leq t_1$.  The analogous statement holds for an attained infimum.
\end{corollary}

\begin{proof}
\sketch
The kernel in Corollary~\ref{cor:euclidean-heat-mean-value-property} is
nonnegative and has total mass one.  Equality of its weighted average with an
attained maximum forces equality throughout every contained heat ball.
Overlapping heat balls and connectedness propagate the equality to every
earlier point.  Apply the same argument to $-u$ for minima.
\end{proof}

\subsection{Strong maximum principles for heat-type systems}

Let $u$ be a nonnegative symmetric bilinear form solving
\[
u_t=\Delta u+\nabla_{X(t)}u+\beta(u,t),
\]
where $\beta$ is locally Lipschitz and satisfies the null-eigenvector
condition.

\begin{proposition}[Strong maximum principle for systems, I]
\label{prop:system-strong-maximum-principle-positive-definite}
\source{chow-et-al-ricci-flow-part-ii:page-0211}
\dcref{ch12:12.47}
\group{Strong System Maximum Principles}

If $u\geq0$ and $u(x_0,0)>0$ at one point, then
$u(x,t)>0$ for every $x$ and every $t>0$.
\end{proposition}

\begin{proof}
\sketch
\uses{thm:manifold-parabolic-strong-maximum-principle}
\source{chow-et-al-ricci-flow-part-ii:page-0211,chow-et-al-ricci-flow-part-ii:page-0212}
Join $x_0$ to an arbitrary $y$ inside a relatively compact connected smooth
domain $\Omega$.  Choose $0\leq\varphi_1\leq\lambda_1(u(\cdot,0))$, positive
at $x_0$ and zero on $\partial\Omega$, and solve
$f_t=\Delta f+Xf-Af$ with initial value $\varphi_1$ and zero boundary value.
The scalar strong maximum principle gives $f>0$ in the interior.  Set
$\widetilde u=u+(\varepsilon e^{At}-f)\widetilde h$.  On the parabolic
boundary $\widetilde u>0$.  Local Lipschitz continuity of $\beta$, the
null-eigenvector condition, and $A$ larger than the Lipschitz constant prevent
a first null eigenvector in the interior by the second derivative test.
Hence $\widetilde u\geq0$.  Letting $\varepsilon\downarrow0$ gives
$u\geq f\widetilde h$, so $u(y,t)>0$.
\end{proof}

\begin{remark}[Locality of the system strong maximum principle]
\label{rem:system-strong-maximum-principle-locality}
\source{chow-et-al-ricci-flow-part-ii:page-0212}
\uses{prop:system-strong-maximum-principle-positive-definite}
\dcref{ch12:12.48}
\group{Strong System Maximum Principles}


The compact connecting domain in the proof shows that this strong maximum
principle is local and requires neither completeness nor global curvature
bounds.
\end{remark}

Let $\lambda_1(u)\leq\cdots\leq\lambda_r(u)$ and put
$\phi_k=\lambda_1+\cdots+\lambda_k$.

\begin{proposition}[Strong maximum principle for systems, II]
\label{prop:system-strong-maximum-principle-eigenvalue-sums}
\source{chow-et-al-ricci-flow-part-ii:page-0213}
\dcref{ch12:12.49}
\group{Strong System Maximum Principles}

If $u\geq0$ and $\phi_k(x_0,0)>0$ for some $k$ and $x_0$, then
$\phi_k(x,t)>0$ for every $x$ and every $t>0$.
\end{proposition}

\begin{proof}
\sketch
\uses{thm:manifold-parabolic-strong-maximum-principle}
\source{chow-et-al-ricci-flow-part-ii:page-0213,chow-et-al-ricci-flow-part-ii:page-0214}
On a relatively compact connecting domain, choose
$0\leq\varphi_k\leq\phi_k(\cdot,0)/k$, positive at $x_0$, and solve the same
scalar Dirichlet problem for $f$.  Put
$\bar u=u+(\varepsilon e^{At}-f)\widetilde h$.  If the sum $\bar\phi_k$ of
its $k$ smallest eigenvalues first vanished, parallel extend an orthonormal
minimizing $k$-frame.  The derivative test, the null-eigenvector condition,
and local Lipschitz continuity yield
\[
0\geq (A-K)k(\varepsilon e^{At}+f),
\]
contradicting $A>K$.  Thus $\bar\phi_k>0$; after
$\varepsilon\downarrow0$, $\phi_k\geq kf>0$ at the arbitrary endpoint.
\end{proof}

\begin{theorem}[Rigidity of the null space]
\label{thm:system-null-space-rigidity}
\source{chow-et-al-ricci-flow-part-ii:page-0215}
\dcref{ch12:12.50}
\group{Strong System Maximum Principles}

Suppose $u\geq0$ solves the system above on a connected manifold.
\begin{enumerate}
\item If $t_2>t_1$, then
\[
\inf_x\operatorname{rank}u(x,t_2)
\geq\sup_x\operatorname{rank}u(x,t_1).
\]
Consequently, for some $\delta>0$, the rank is constant in space and time on
$M\times(0,\delta)$.
\item On this interval $\operatorname{null}u(t)$ is a smooth subbundle,
smoothly depending on time.
\item It is invariant under spatial parallel translation, constant in time,
and
\[
\operatorname{null}u(t)\subset
\operatorname{null}\beta(u(t),t).
\]
\end{enumerate}
\end{theorem}

\begin{proof}
\sketch
\uses{prop:system-strong-maximum-principle-eigenvalue-sums}
\source{chow-et-al-ricci-flow-part-ii:page-0216,chow-et-al-ricci-flow-part-ii:page-0217,chow-et-al-ricci-flow-part-ii:page-0218,chow-et-al-ricci-flow-part-ii:page-0219}
If $\ell=\sup_x\operatorname{rank}u(x,t_1)$, then
$\phi_{r-\ell+1}>0$ somewhere at $t_1$; Proposition~\ref{prop:system-strong-maximum-principle-eigenvalue-sums}
makes it positive everywhere at $t_2$, proving the rank inequality.  Its
integer-valued sides imply short-time constancy.  The positive and zero
spectral spaces are then separated by a uniform local spectral gap, so
spectral projection makes the null space a smooth space-time subbundle.

Let $W$ be a local smooth null section.  Differentiating $u(W)=0$ and using
the evolution equation gives
\[
2g^{ij}u(\nabla_iW,\nabla_jW)+\beta(u,t)(W,W)=0.
\]
Both terms are nonnegative, hence $\nabla_YW$ is null for every $Y$ and
$\beta(u,t)(W,W)=0$.  Solving the linear coefficient ODE along any curve
shows invariance under parallel translation.  This invariance implies
$u(\nabla_jW)=0$ and, after differentiating once more,
$(\Delta u)(W)=0$.  The evolution equation now gives
$u(\partial_tW)=0$.  Solving the analogous coefficient ODE in time proves
that the null subbundle is constant in time.
\end{proof}

\begin{remark}[Infinitesimal parallel invariance]
\label{rem:null-space-infinitesimal-parallel-invariance}
\source{chow-et-al-ricci-flow-part-ii:page-0215}
\uses{thm:system-null-space-rigidity}
\dcref{ch12:12.51}
\group{Strong System Maximum Principles}


For every smooth local null section $W$ and tangent vector $Y$,
$\nabla_YW$ is again a null section.
\end{remark}

\section{Applications to the curvature operator under Ricci flow}

\begin{remark}[Curvature consequences on noncompact solutions]
\label{rem:noncompact-ricci-flow-curvature-consequences}
\source{chow-et-al-ricci-flow-part-ii:page-0220}
\uses{thm:noncompact-two-tensor-maximum-principle, thm:weinberger-hamilton-noncompact-systems}
\dcref{ch12:12.52}
\group{Ricci Flow Applications}


For complete bounded-curvature Ricci flows on noncompact manifolds, the weak
system maximum principles preserve lower scalar-curvature bounds,
nonnegative curvature operator, and in dimension three nonnegative Ricci
curvature and the Hamilton--Ivey pinching inequality.  They likewise preserve
nonnegative bisectional curvature under K\"ahler--Ricci flow.
\end{remark}

\begin{theorem}[Strong maximum principle for the curvature operator]
\label{thm:curvature-operator-strong-maximum-principle}
\source{chow-et-al-ricci-flow-part-ii:page-0220,chow-et-al-ricci-flow-part-ii:page-0221}
\dcref{ch12:12.53}
\group{Ricci Flow Applications}

Let $g(t)$ be a Ricci flow on a connected manifold without boundary and assume
$\Rm(g(t))\geq0$.
\begin{enumerate}
\item If $\Rm(x_0,0)>0$ at one point, then $\Rm(x,t)>0$ everywhere for every
$t>0$.
\item For some $\delta>0$, $\operatorname{image}\Rm(g(t))\subset
\Lambda^2T^*M$ is a smooth subbundle on $(0,\delta)$, invariant under parallel
translation and constant in time.  Each fiber is a Lie subalgebra of
$\Lambda^2T_x^*M$.
\end{enumerate}
\end{theorem}

\begin{proof}
\sketch
\uses{prop:system-strong-maximum-principle-positive-definite,thm:system-null-space-rigidity}
\source{chow-et-al-ricci-flow-part-ii:page-0221}
After Uhlenbeck identification, the curvature operator satisfies a system
whose reaction term obeys the null-eigenvector condition.  The first assertion
is Proposition~\ref{prop:system-strong-maximum-principle-positive-definite};
the smooth, parallel, time-independent image follows from
Theorem~\ref{thm:system-null-space-rigidity} applied to the orthogonal null
space.

For the Lie assertion, diagonalize
$\Rm=\sum_\alpha\lambda_\alpha\varphi^\alpha\otimes\varphi^\alpha$ with
$\lambda_\alpha=0$ for $\alpha\leq p$ and positive otherwise.  Null-space
rigidity and the reaction equation imply
$\operatorname{null}\Rm\subset\operatorname{null}\Rm^\#$.  Thus, for
$\alpha\leq p$,
\[
0=(\Rm^\#)_{\alpha\alpha}
=\sum_{\gamma,\delta}(C_\alpha^{\gamma\delta})^2
\lambda_\gamma\lambda_\delta,
\]
so $C_\alpha^{\gamma\delta}=0$ whenever $\gamma,\delta>p$.  Therefore the
bracket of two vectors in
$\operatorname{span}\{\varphi^\beta:\beta>p\}=\operatorname{image}\Rm$
again lies in that span.
\end{proof}

\begin{remark}[Time interval for curvature-operator rigidity]
\label{rem:curvature-operator-image-time-interval}
\source{chow-et-al-ricci-flow-part-ii:page-0221}
\uses{thm:curvature-operator-strong-maximum-principle}
\dcref{ch12:12.54}
\group{Ricci Flow Applications}


For complete bounded-curvature flows, nonnegativity at time zero supplies the
hypothesis $\Rm(g(t))\geq0$.  The image subbundle in the theorem need not be
constant on the whole interval $(0,T)$.
\end{remark}

\begin{remark}[Positive bisectional curvature]
\label{rem:kahler-ricci-positive-bisectional-curvature}
\source{chow-et-al-ricci-flow-part-ii:page-0222}
\uses{prop:system-strong-maximum-principle-positive-definite}
\dcref{ch12:12.55}
\group{Ricci Flow Applications}


For a K\"ahler--Ricci flow on a connected manifold without boundary,
nonnegative bisectional curvature everywhere and positive bisectional
curvature somewhere at time zero imply positive bisectional curvature
everywhere for every positive time.
\end{remark}

\begin{example}[Exercise: backward-time proof for continuously varying sets]
\label{ex:backward-time-convex-set-maximum-principle}
\dcref{ch12:12.37}
\group{Weak System Maximum Principles}
\source{chow-et-al-ricci-flow-part-ii:page-0204}
If $\mathcal K(t)$ varies continuously, $X=0$, and $F$ is independent of time,
Theorem~\ref{thm:weinberger-hamilton-noncompact-systems} also follows by
estimating the backward upper derivative of the distance function.
\end{example}
\begin{proof}
\uses{lem:ricci-flow-barrier-supersolution}
\source{chow-et-al-ricci-flow-part-ii:page-0204,chow-et-al-ricci-flow-part-ii:page-0205,chow-et-al-ricci-flow-part-ii:page-0206}
At a point $(x_0,t_0)$ maximizing
$\rho_{t_0}(u)-\varepsilon\phi$, compare with time $t_0-h$.  Convexity of
$\rho_{t_0-h}$ and the second derivative test imply that the contribution of
$-h\Delta u$ is at least $-h\varepsilon\Delta\phi+o(h)$.  If $v_h$ is the
nearest point in $\mathcal K(t_0-h)_{x_0}$, the ODE solution from $v_h$ lies
in $\mathcal K(t_0)_{x_0}$ and equals $v_h+hF(v_h)+o(h)$.  The Lipschitz
property of $F$ therefore gives
\[
\rho_{t_0}(u)-\rho_{t_0-h}(u-hF(u))
\leq hC\rho_{t_0}(u)+o(h).
\]
After the barrier terms are restored, $d^+s/dt\leq Cs$; since $s(0)=0$, one
gets $s\leq0$, and then lets $\varepsilon\downarrow0$.
\end{proof}

\begin{example}[Exercise: constructing the first barrier]
\label{ex:construct-general-metric-barrier}
\dcref{ch12:12.8}
\group{Barrier Functions}
\uses{lem:general-metric-barrier-supersolution}
\source{chow-et-al-ricci-flow-part-ii:page-0172}
The function $\phi=e^{Bt+\gamma h}$ in the preceding proof satisfies the
claimed differential and growth inequalities after the indicated choices of
$\gamma$ and $B$.
\end{example}

\begin{example}[Exercise: the $L^p$ energy inequality]
\label{ex:derive-lp-heat-energy-inequality}
\dcref{ch12:12.27}
\group{Heat Equation Uniqueness}
\uses{thm:lp-heat-subsolution-vanishing}
\source{chow-et-al-ricci-flow-part-ii:page-0184}
The energy inequality in the proof of
Theorem~\ref{thm:lp-heat-subsolution-vanishing} follows from
\[
-4\int_Mu^{p/2}\varphi_r\langle\nabla\varphi_r,
\nabla u^{p/2}\rangle
\leq\frac{2(p-1)}p\int_M\varphi_r^2|\nabla u^{p/2}|^2
+\frac{2p}{p-1}\int_Mu^p|\nabla\varphi_r|^2.
\]
\end{example}
